% 附录A 第二次量子化
\chapter{第二次量子化}

这里复习正文中通用的二次量子化的标准定义与基本关系。

\section{Fock 态}

给定正交归一的单粒子基底 $\{|\phi_{i}\rangle\}_{i=1}^{N}$：

\begin{equation*}
\langle \phi_ {i} | \phi_ {j} \rangle = \delta_ {ij}.
\tag{A.1}
\end{equation*}

态 $i$ 的产生算符是 $a_{i}^{\dagger}$，其厄米共轭是湮灭算符 $a_{i}$。二者都相对真空态 $|0\rangle_{i}$ 定义：

\begin{equation*}
\left| \phi_ {i} \right\rangle = a _ {i} ^ {\dagger} \left| 0 \right\rangle_{i} \, , \quad a _ {i} \left| 0 \right\rangle_{i} = 0.
\tag{A.2}
\end{equation*}

数算符定义为

\begin{equation*}
n _ {i} \equiv a _ {i} ^ {\dagger} a _ {i}.
\tag{A.3}
\end{equation*}

利用对易规则

\begin{equation*}
\left[ a _ {i} ^ {\dagger} a _ {i}, \left( a _ {i} ^ {\dagger} \right) ^ {n} \right] = n \left( a _ {i} ^ {\dagger} \right) ^ {n} ,
\tag{A.4}
\end{equation*}

可以验证反复作用 $a_{i}^{\dagger}$（$a_{i}$）增加（减少）态 $\phi_{i}$ 中的粒子数。以占据数 $n_{i}$ 标记的 Fock 态定义为

\begin{equation*}
\left| n _ {1}, n _ {2}, \dots \right\rangle = \prod_ {i} \frac {\left( a _ {i} ^ {\dagger} \right) ^ {n _ {i}}}{\sqrt {n _ {i} !}} \left| 0 \right\rangle_{i}.
\tag{A.5}
\end{equation*}

对多玻色子（或费米子），Fock 态的对称（反对称）性由对易（反对易）关系强制：

\begin{equation*}
\left[ a _ {i}, a _ {j} ^ {\dagger} \right]_{\eta} = a _ {i} a _ {j} ^ {\dagger} - \eta \, a _ {j} ^ {\dagger} a _ {i} = \delta_{ij} ,
\tag{A.6}
\end{equation*}

\begin{equation*}
\left[ a _ {i}, a _ {j} \right]_{\eta} = 0 ,
\tag{A.7}
\end{equation*}

其中

\begin{equation*}
\eta = \left\{ \begin{array}{ll} 1 & \text{玻色子} \\ - 1 & \text{费米子} \end{array} \right.
\tag{A.8}
\end{equation*}

占据数受总粒子数约束限制：

\begin{equation*}
\sum_ {i} n _ {i} = N.
\tag{A.9}
\end{equation*}

对易（反对易）关系定义了二次量子化算符的一个代数。该代数在二次量子化算符的正则变换下不变。这里只考虑正规（normal）变换\footnote{关于 Bogoliubov 变换的专著见文献注释。}。考虑单粒子基底上的幺正变换 $U$：

\begin{equation*}
a _ {\alpha} ^ {\dagger} = \sum_ {i} U _ {\alpha i} a _ {i} ^ {\dagger} \, , \quad U ^ {\dagger} U = I.
\tag{A.10}
\end{equation*}

容易验证 $\{a_{\alpha}^{\dagger}\}$ 与 $\{a_{\alpha}\}$ 满足正则代数 (A.6) 与 (A.7)。

\section{正规双线性算符}

双线性算符由

\begin{equation*}
\hat {A} = \sum_ {ij} a _ {i} ^ {\dagger} A _ {ij} a _ {j} \equiv \mathbf {a} ^ {\dagger} \cdot A \cdot \mathbf {a}
\tag{A.11}
\end{equation*}

给出，$A$ 是基底 $\phi_{i}$ 中的厄米矩阵，即单粒子算符。

玻色子或费米子的双线性算符与线性算符的对易关系特别简单。矢量 $\mathbf{v}$ 定义线性算符 $\hat{v}^{\dagger}$：

\begin{equation*}
\hat {\mathbf {v}} ^ {\dagger} = \sum_ {i} v _ {i} a _ {i} ^ {\dagger} = \mathbf {v} \cdot \mathbf {a} ^ {\dagger}.
\tag{A.12}
\end{equation*}

把 (A.11) 与 (A.12) 对易：

\begin{equation*}
\left[ \hat {A}, \hat {\mathbf {v}} ^ {\dagger} \right] = \left( A \mathbf {v} \right) \cdot \mathbf {a} ^ {\dagger}.
\tag{A.13}
\end{equation*}

若 $\mathbf{v}$ 是 $A$ 的本征值为 $v$ 的本征矢量，则 $\hat{v}^{\dagger}$ 是 $[A,\cdot]$ 的本征值为 $v$ 的本征算符。这可用于计算 $\hat{v}^{\dagger}$ 在转动下的变换：

\begin{align*}
e ^ {i\theta \hat {A}} \hat {\mathbf {v}} ^ {\dagger} e ^ {-i\theta \hat {A}} & = \hat {\mathbf {v}} ^ {\dagger} + i\theta \left[ \hat {A}, \hat {\mathbf {v}} ^ {\dagger} \right] + \frac {( i\theta ) ^ {2}}{2} \left[ \hat {A}, \left[ \hat {A}, \hat {\mathbf {v}} ^ {\dagger} \right] \right] + \dots\\
& = e ^ {iv\theta} \hat {\mathbf {v}} ^ {\dagger} .
\tag{A.14}
\end{align*}

幺正矩阵 $U$ 可用一组厄米生成元 $A_{\alpha}$ 与参数 $\theta_{\alpha}$ 写出：

\begin{equation*}
U _ {\theta} = e ^ {i \sum_ {\alpha} \theta_{\alpha} A_{\alpha}}.
\tag{A.15}
\end{equation*}

相应的多粒子算符 $\hat{U}$ 定义为

\begin{equation*}
\hat {U} _ {\theta} = e ^ {i \sum_ {\alpha} \theta_{\alpha} \hat {A}_{\alpha}}.
\tag{A.16}
\end{equation*}

利用 (A.14)，容易证明 $\hat{\mathbf{v}}^{\dagger}$ 在 $\hat{U}$ 下按简单方式变换：

\begin{equation*}
\hat {U} _ {\theta} \hat {\mathbf {v}} ^ {\dagger} \hat {U} _ {\theta} ^ {- 1} = \left( U _ {\theta} \mathbf {v} \right) \cdot \mathbf {a} ^ {\dagger}.
\tag{A.17}
\end{equation*}

两个双线性算符之间的对易关系是

\begin{equation*}
\left[ \hat {A}, \hat {B} \right] = \mathbf {a} ^ {\dagger} [ A, B ] \mathbf {a}.
\tag{A.18}
\end{equation*}

(A.17) 与 (A.18) 对费米子与玻色子都成立。

\subsection*{例：自旋代数}

双线性自旋算符是

\begin{equation*}
S ^ {\alpha} = \frac {1}{2} \sum_{ss' = 1} ^ {2} a _ {s} ^ {\dagger} \sigma_{ss'} ^ {\alpha} a _ {s '} \, , \quad \alpha = x, y, z ,
\tag{A.19}
\end{equation*}

其中

\begin{equation*}
\left( \sigma^ {x}, \sigma^ {y}, \sigma^ {z} \right) = \left[ \left( \begin{array}{cc} 0 & 1 \\ 1 & 0 \end{array} \right) , \left( \begin{array}{cc} 0 & -i \\ i & 0 \end{array} \right) , \left( \begin{array}{cc} 1 & 0 \\ 0 & -1 \end{array} \right) \right].
\tag{A.20}
\end{equation*}

计算 Pauli 矩阵的对易关系并用 (A.18)，可以验证 $S^{\alpha}$ 满足角动量关系

\begin{equation*}
\left[ S ^ {\alpha}, S ^ {\beta} \right] = i \epsilon^ {\alpha \beta \gamma} S ^ {\gamma} ,
\tag{A.21}
\end{equation*}

$\epsilon$ 是全反对称张量。

\section{无相互作用哈密顿量}

正规二次型哈密顿量为

\begin{equation*}
\mathcal {H} = \sum_ {ii '} a _ {i} ^ {\dagger} H _ {ii '} a _ {i '} = \sum_ {k} \left( \epsilon_ {k} - \mu \right) a _ {k} ^ {\dagger} a _ {k} ,
\tag{A.22}
\end{equation*}

$H$ 厄米，$\epsilon_{k}, a_{k}$ 分别是它的能量与本征算符，$\mu$ 是化学势。配分函数为

\begin{align*}
Z & = \operatorname{Tr} e ^ {- \mathcal {H} /T} = \prod_ {k} \sum_{n _ {k} = 0} ^ {n_{\mathrm{max}}} \exp \left[ - \frac {\left( \epsilon_ {k} - \mu \right) n _ {k}}{T} \right]\\
& = \prod_ {k} \left( 1 - \eta \, e ^ {- \left( \epsilon_ {k} - \mu \right) /T} \right) ^ {- \eta} .
\tag{A.23}
\end{align*}

(A.23) 中 $n_{\mathrm{max}}$ 对玻色子为无穷，对费米子为 1。

由 (A.23)，自由能为

\begin{equation*}
F \equiv -T \ln Z = \eta T \sum_ {k} \ln \left( 1 - \eta \, e ^ {- \left( \epsilon_ {k} - \mu \right) /T} \right).
\tag{A.24}
\end{equation*}

平衡占据概率为

\begin{align*}
\langle n _ {k} \rangle & = \frac {1}{Z} \operatorname{Tr} \left[ a _ {k} ^ {\dagger} a _ {k} e ^ {- H /T} \right] = \frac {\partial F}{\partial \epsilon_ {k}}\\
& = \left( e ^ {\left( \epsilon_ {k} - \mu \right) /T} - \eta \right) ^ {- 1} ,
\tag{A.25}
\end{align*}

$\eta = \pm1$ 时分别是 Bose 函数与 Fermi 函数。

\begin{exercises}

\item 两个 $N$ 费米子 Fock 态 $\Phi_1 = |\{n_i\} \rangle$ 与 $\Phi_2 = |\{n_\alpha\} \rangle$ 分别相对基底 $\{\phi_i\}$ 与 $\{\phi_\alpha\}$ 定义，两基底由幺正矩阵 $U$ 联系：

\begin{equation*}
\phi_ {\alpha} = \sum_ {i} U _ {\alpha i} \phi _ {i}.
\tag{A.26}
\end{equation*}

用矩阵 $U$ 表出交叠 $\langle\Phi_{1}|\Phi_{2}\rangle$。

\item 自旋转动算符为

\begin{equation*}
R = \exp \left[ iS ^ {z} \phi \right] \exp \left[ iS ^ {y} \theta \right] \exp \left[ iS ^ {z} \chi \right].
\tag{A.27}
\end{equation*}

注意 $R$ 是定义在费米子 Fock（占据数）空间中的算符。用 (A.17)，求 $SU(2)$ 变换矩阵 $U_{ss'}$ 的显式表达式：

\begin{equation*}
\left( c ' \right)_{s} = \left( R c R ^ {- 1} \right)_{s} = \sum_ {s '} U _ {ss'} \left( \phi, \theta, \chi \right) c _ {s '}.
\tag{A.28}
\end{equation*}

\item 求双线性形式

\begin{equation*}
c _ {1 \uparrow} ^ {\dagger} c _ {2 \downarrow} ^ {\dagger} - c _ {1 \downarrow} ^ {\dagger} c _ {2 \uparrow} ^ {\dagger}
\tag{A.29}
\end{equation*}

在整体自旋转动下的变换，其中 $c_{1s}$ 与 $c_{2s}$ 代表两个不同轨道上的电子。

\item 证明：对 (A.28) 给出的任何整体转动，顺磁 Fermi 气体态保持不变：

\begin{equation*}
\prod_{s, \, | k | < \pi / 2} \left( c _ {ks} ^ {\dagger} \right) ' \left| 0 \right\rangle = \prod_{s, \, | k | < \pi / 2} c _ {ks} ^ {\dagger} \left| 0 \right\rangle .
\tag{A.30}
\end{equation*}

\end{exercises}

\begin{chapbib}
\item 广义 Bogoliubov 变换的讨论见
  R. Balian and E. Brézin, Nuovo Cimento B \textbf{LXIV}, 37 (1969)。
\end{chapbib}
